% --------------------------------------------------
\subsubsection{Тогтмол корреляцын зорилтот матриц}
% --------------------------------------------------

Тогтмол корреляцын зорилтот матриц нь хөрөнгө тус бүрийн
өгөөжийн түүврийн вариацыг хадгалж, харин хөрөнгүүдийн
хоорондын корреляцыг тэдгээрийн дундаж утгаар
төлөөлүүлэн байгуулдаг зорилтот матриц юм
\citep{ledoitwolf2004b}.

$i$ болон $j$ дугаар хөрөнгийн өгөөжийн
түүврийн корреляцын коэффициентийг
\[
r_{ij}
=
\frac{s_{ij}}{\sqrt{s_{ii}s_{jj}}}
\]
гэж тодорхойлно \citep{ledoitwolf2004b}.

Хөрөнгүүдийн хоорондын түүврийн
корреляцуудын дундажийг
\[
\bar r
=
\frac{2}{N(N-1)}
\sum_{i=1}^{N-1}
\sum_{j=i+1}^{N}r_{ij}
\]
гэж тооцно.

Тогтмол корреляцын зорилтот матрицын
элементүүдийг
\[
(T_{\mathrm{corr}})_{ij}
=
t_{ij}
=
\begin{cases}
	s_{ii}, & i=j,\\
	\bar r\sqrt{s_{ii}s_{jj}}, & i\neq j
\end{cases}
\]
гэж тодорхойлно.

Иймээс уг зорилтот матриц нь хөрөнгө тус бүрийн
өгөөжийн түүврийн вариацыг өөрчлөхгүй хадгалж,
харин хөрөнгүүдийн хоорондын корреляцыг
нэг ижил дундаж утгаар орлуулна.

Тогтмол корреляцын зорилтот матрицад
суурилсан шугаман агшилтын үнэлэгчийг
\[
\widehat{\Sigma}_{\mathrm{corr}}
=
(1-\widehat{\alpha}_{\mathrm{corr}})S
+
\widehat{\alpha}_{\mathrm{corr}}T_{\mathrm{corr}},
\qquad
0\leq\widehat{\alpha}_{\mathrm{corr}}\leq1
\]
гэж тодорхойлно. Энд
$\widehat{\alpha}_{\mathrm{corr}}$ нь тогтмол
корреляцын зорилтот матрицад харгалзах
агшилтын эрчим юм.

\paragraph{Агшилтын эрчмийн үнэлгээ.}

Тогтмол корреляцын зорилтот матриц нь
түүврийн ковариацын матрицын элементүүдээс
хамааран байгуулагддаг тул агшилтын
эрчмийг үнэлэхдээ түүврийн ковариацын
матриц болон зорилтот матрицын
хам хэлбэлзлийг харгалзан үзнэ.

Түүврийн ковариацын матриц болон
тогтмол корреляцын зорилтот матрицын
хоорондын нормчилсон квадрат
Фробениусын зайг
\[
\widehat d_{\mathrm{corr}}
=
\|S-T_{\mathrm{corr}}\|^2
=
\frac{1}{N}
\sum_{i=1}^{N}\sum_{j=1}^{N}
(s_{ij}-t_{ij})^2
=
\frac{1}{N}
\sum_{\substack{i,j=1\\i\neq j}}^{N}
\left(
s_{ij}-\bar r\sqrt{s_{ii}s_{jj}}
\right)^2
\]
гэж тодорхойлно.

Түүврийн ковариацын матрицын
хүлээгдэж буй квадрат алдааг өмнөх хэсгүүдэд үзсэнтэй адилаар 
\[
\widehat b_{\mathrm{corr}}
=
\frac{1}{nN}
\sum_{i=1}^{N}\sum_{j=1}^{N}
\left[
\frac{1}{n}
\sum_{t=1}^{n}
\left\{
(y_{ti}-\overline{Y}_i)
(y_{tj}-\overline{Y}_j)-s_{ij}
\right\}^2
\right]
\]
гэж бичигдэнэ.

Тогтмол корреляцын зорилтот матриц
нь түүврийн ковариацын матрицаас
хамаардаг учраас тэдгээрийн
харгалзах элементүүдийн
асимптотик ковариацыг тооцох
шаардлагатай.

Энэ хэмжигдэхүүний нормчилсон
үнэлгээг $\widehat c_{\mathrm{corr}}$
гэж тэмдэглэе.

Эхлээд
\[
\widehat{\vartheta}_{ii,ij}
=
\frac{1}{n}
\sum_{t=1}^{n}
\left\{
(y_{ti}-\overline{Y}_i)^2-s_{ii}
\right\}
\left\{
(y_{ti}-\overline{Y}_i)
(y_{tj}-\overline{Y}_j)-s_{ij}
\right\}
\]
болон
\[
\widehat{\vartheta}_{jj,ij}
=
\frac{1}{n}
\sum_{t=1}^{n}
\left\{
(y_{tj}-\overline{Y}_j)^2-s_{jj}
\right\}
\left\{
(y_{ti}-\overline{Y}_i)
(y_{tj}-\overline{Y}_j)-s_{ij}
\right\}
\]
гэж тодорхойлно.

Эдгээр нь түүврийн вариац болон
ковариацын элементүүдийн
асимптотик ковариацыг үнэлэхэд
ашиглагдана.

Ledoit болон Wolf (2004) нь түүврийн ковариацын
матриц болон тогтмол корреляцын зорилтот матрицын
асимптотик хам хэлбэлзлийг тооцох аргыг
боловсруулсан \citep{ledoitwolf2004b}.
Тэдгээрийн Appendix B-д өгсөн үр дүнд тулгуурлан
\[
	\widehat c_{\mathrm{corr}}
	=
	\frac{1}{nN}
	\Bigg[
	\sum_{i=1}^{N}
	\widehat{\operatorname{Var}}
	\left((y_{1i}-\overline{Y}_i)^2\right)+
	\frac{\bar r}{2}
	\sum_{\substack{i,j=1\\i\neq j}}^{N}
	\left(
	\sqrt{\frac{s_{jj}}{s_{ii}}}
	\widehat{\vartheta}_{ii,ij}
	+
	\sqrt{\frac{s_{ii}}{s_{jj}}}
	\widehat{\vartheta}_{jj,ij}
	\right)
	\Bigg]
\]
гэж үнэлнэ.

Энд
\[
\widehat{\operatorname{Var}}
\left((y_{1i}-\overline{Y}_i)^2\right)
=
\frac{1}{n}
\sum_{t=1}^{n}
\left\{
(y_{ti}-\overline{Y}_i)^2-s_{ii}
\right\}^2
\]
юм.

Иймээс тогтмол корреляцын
зорилтот матрицад харгалзах
агшилтын эрчмийн практик үнэлгээг
\[
\widehat{\alpha}_{\mathrm{corr}}
=
\max\left\{
0,\,
\min\left\{
1,\,
\frac{
	\widehat b_{\mathrm{corr}}
	-
	\widehat c_{\mathrm{corr}}
}{
	\widehat d_{\mathrm{corr}}
}
\right\}
\right\}
\]
гэж байгуулна \citep{ledoitwolf2004b}.

Энд $\widehat b_{\mathrm{corr}}$ нь
түүврийн ковариацын матрицын
үнэлгээний квадрат алдааны
асимптотик ойролцоо үнэлгээ,
$\widehat c_{\mathrm{corr}}$ нь
түүврийн ковариацын матриц болон
зорилтот матрицын харгалзах
элементүүдийн асимптотик
ковариацуудын нийлбэрийн
нормчилсон үнэлгээ,
$\widehat d_{\mathrm{corr}}$ нь
хоёр матрицын хоорондын
нормчилсон квадрат Фробениусын
зай юм.

Эцэст нь агшилтын ковариацын
матрицыг
\[
\widehat{\Sigma}_{\mathrm{corr}}
=
(1-\widehat{\alpha}_{\mathrm{corr}})S
+
\widehat{\alpha}_{\mathrm{corr}}
T_{\mathrm{corr}}
\]
гэж тооцно.

\paragraph{Тооцоолох алгоритм.}

Тогтмол корреляцын зорилтот матрицад
суурилсан агшилтын ковариацын матрицыг
дараах дарааллаар тооцоолно.

\begin{enumerate}
	\item Тогтмол корреляцын зорилтот
	матрицыг байгуулах.
	
	Өмнөх хэсэгт тооцсон түүврийн
	ковариацын матриц $S=(s_{ij})$-ийг
	ашиглан
	\[
	r_{ij}
	=
	\frac{s_{ij}}{\sqrt{s_{ii}s_{jj}}},
	\qquad
	\bar r
	=
	\frac{2}{N(N-1)}
	\sum_{i=1}^{N-1}
	\sum_{j=i+1}^{N}r_{ij}
	\]
	гэж тооцно.
	
	Дараа нь
	\[
	t_{ij}
	=
	\begin{cases}
		s_{ii}, & i=j,\\
		\bar r\sqrt{s_{ii}s_{jj}}, & i\neq j
	\end{cases}
	\]
	гэж тодорхойлж,
	$T_{\mathrm{corr}}=(t_{ij})$
	матрицыг байгуулна.
	
	\item Ковариацын матрицын
	квадрат алдааг үнэлэх.
	
	\[
	\widehat b_{\mathrm{corr}}
	=
	\frac{1}{nN}
	\sum_{i=1}^{N}\sum_{j=1}^{N}
	\left[
	\frac{1}{n}
	\sum_{t=1}^{n}
	\left\{
	(y_{ti}-\overline{Y}_i)
	(y_{tj}-\overline{Y}_j)-s_{ij}
	\right\}^2
	\right]
	\]
	гэж тооцно.
	
	\item Түүврийн ковариацын матриц
	болон зорилтот матрицын
	хам хэлбэлзлийг үнэлэх.
	
	\[
	\widehat{\vartheta}_{ii,ij}
	=
	\frac{1}{n}
	\sum_{t=1}^{n}
	\left\{
	(y_{ti}-\overline{Y}_i)^2-s_{ii}
	\right\}
	\left\{
	(y_{ti}-\overline{Y}_i)
	(y_{tj}-\overline{Y}_j)-s_{ij}
	\right\}
	\]
	болон
	\[
	\widehat{\vartheta}_{jj,ij}
	=
	\frac{1}{n}
	\sum_{t=1}^{n}
	\left\{
	(y_{tj}-\overline{Y}_j)^2-s_{jj}
	\right\}
	\left\{
	(y_{ti}-\overline{Y}_i)
	(y_{tj}-\overline{Y}_j)-s_{ij}
	\right\}
	\]
	гэж тооцно.
	
	Дараа нь
	\[
		\widehat c_{\mathrm{corr}}
		=
		\frac{1}{nN}
		\Bigg[
		\sum_{i=1}^{N}
		\widehat{\operatorname{Var}}
		\left((y_{1i}-\overline{Y}_i)^2\right)+
		\frac{\bar r}{2}
		\sum_{\substack{i,j=1\\i\neq j}}^{N}
		\left(
		\sqrt{\frac{s_{jj}}{s_{ii}}}
		\widehat{\vartheta}_{ii,ij}
		+
		\sqrt{\frac{s_{ii}}{s_{jj}}}
		\widehat{\vartheta}_{jj,ij}
		\right)
		\Bigg]
	\]
	гэж үнэлнэ.
	
	\item Зорилтот матрицаас хазайх
	квадрат зайг тооцох.
	
	\[
	\widehat d_{\mathrm{corr}}
	=
	\frac{1}{N}
	\sum_{\substack{i,j=1\\i\neq j}}^{N}
	(s_{ij}-t_{ij})^2.
	\]
	
	\item Агшилтын эрчмийг үнэлэх.
	
	$\widehat d_{\mathrm{corr}}>0$
	үед
	\[
	\widehat{\alpha}_{\mathrm{corr}}
	=
	\max\left\{
	0,\,
	\min\left\{
	1,\,
	\frac{
		\widehat b_{\mathrm{corr}}
		-
		\widehat c_{\mathrm{corr}}
	}{
		\widehat d_{\mathrm{corr}}
	}
	\right\}
	\right\}.
	\]
	
	\item Агшилтын ковариацын
	матрицыг байгуулах.
	
	\[
	\widehat{\Sigma}_{\mathrm{corr}}
	=
	(1-\widehat{\alpha}_{\mathrm{corr}})S
	+
	\widehat{\alpha}_{\mathrm{corr}}
	T_{\mathrm{corr}}.
	\]
\end{enumerate}